\chapter{\ResearchThreeTitle}
\label{rq3:chapter}

前两章分别研究了传输资源分配和语义内容选择。本章针对数据丢失、确认延迟或丢失造成的交付不确定性，在候选内容给定时兼顾任务按时正确完成率与通信能耗。发送端依据发送记录与累计确认，决定新发、重传、等待或停止，系统结构如图~\ref{rq3:fig:system} 所示。

\begin{figure}[htbp]
\centering
% Functional view: candidate descriptions enter the controller; selected
% messages are scheduled over the data link and cumulative feedback returns.
\begin{tikzpicture}[x=1mm,y=1mm,
  font=\fontsize{9bp}{11.5bp}\selectfont,
  box/.style={draw,line width=.45pt,fill=white,align=center,inner sep=1mm},
  flow/.style={-{Stealth[length=1.5mm,width=1.2mm]},line width=.55pt},
  ack/.style={flow,dashed},
  edgeword/.style={fill=white,inner sep=.3mm}]
  \path[use as bounding box] (0,0) rectangle (148,-52);
  \node[box,text width=31mm,minimum height=10mm] (a) at (18.5,-14)
    {感知节点 A\\局部语义证据};
  \node[box,text width=31mm,minimum height=10mm] (b) at (18.5,-27)
    {感知节点 B\\局部语义证据};
  \node[box,text width=31mm,minimum height=10mm] (c) at (18.5,-40)
    {其他感知节点\\局部语义证据};
  \node[box,text width=42mm,minimum height=22mm] (ctrl) at (75,-27)
    {\textbf{C-BFC 控制器}\\接收信念与集合效用\\能量价格与两步前瞻\\新发/重传/等待/停止};
  \node[box,text width=34mm,minimum height=20mm] (rx) at (129,-27)
    {接收中心\\唯一证据融合\\任务分类结果};
  \draw[flow] (a.east) -- (43,-14) |- ($(ctrl.west)+(0,6mm)$);
  \draw[flow] (b.east) -- node[edgeword,above] {候选描述} (ctrl.west);
  \draw[flow] (c.east) -- (43,-40) |- ($(ctrl.west)+(0,-6mm)$);
  \draw[flow] (ctrl.east) -- node[edgeword,above] {调度发送} (rx.west);
  \draw[ack] (rx.north) -- (129,-4.5) -- node[edgeword,above] {延迟或丢失的累计确认} (75,-4.5) -- (ctrl.north);
  \node[align=center] at (74,-49.5) {节点发送能量与任务截止时间约束};
\end{tikzpicture}

\bicaption{语义可靠交付的逻辑结构与反馈关系}{Logical architecture and feedback flow of reliable semantic delivery}
\label{rq3:fig:system}
\end{figure}

针对接收集合不可直接观测、重复传输消耗有限电量的问题，本章设计一致性感知的信念反馈控制方法 C-BFC（Consistency-aware Belief Feedback Control）。该方法估计不同接收集合的任务效用，依据发送和反馈记录更新接收概率，并用有限前瞻选择后续动作。虚拟能量队列记录节点在已完成任务中的支出，使当前交付决策考虑后续任务的通信需要。

\section{系统模型与问题描述}
\label{rq3:sec:model}

\subsection{观测任务与传输反馈}
\label{rq3:subsec:tasks}
\label{rq3:subsec:channel}
\label{rq3:subsec:feedback}

考虑由感知节点集合 $\mathcal N=\{1,\ldots,N\}$ 与地面接收中心 $d$ 构成的协同观测系统，依次处理 $N_{\mathrm{task}}$ 个连续任务，其中第 $k$ 个任务的相对截止时间为 $D_k$。各节点的通信电量跨任务连续扣减，而接收状态与确认记录则在每个新任务开始时重置。

对于任务 $k$，上层内容选择模块依据各节点的局部观测与语义粒度，预先生成候选证据集合
\begin{equation}
  \mathcal{M}_k=\{m_{k,1},m_{k,2},\ldots,m_{k,K_k}\}.
  \label{rq3:eq:candidates}
\end{equation}
每条候选证据 $m$ 包含发送节点 $s(m)$、词元长度 $\ell_m$、表达粒度（标签、简洁声明、摘要或完整证据）以及对应的信念更新量。在交付控制阶段，候选内容、融合次序与融合规则均保持固定，控制器只需决定发送哪条候选，从而隐式确定对应的发送节点与传输开销。

记当前任务内决策时刻 $t$ 接收中心实际收到的唯一证据集合为 $S_{k,t}\subseteq\mathcal{M}_k$，则可能的接收状态空间为
\begin{equation}
  \states_k=2^{\mathcal{M}_k},\qquad |\states_k|=2^{K_k}.
  \label{rq3:eq:state-space}
\end{equation}
接收中心按既定次序对 $S_{k,t}$ 中的唯一证据执行融合，并输出预测类别 $c_k(S_{k,t})$。由于同一条证据仅融合一次，最终分类结果完全取决于接收集合本身，不受重复到达次数或到达先后顺序的影响。记该任务的参考类别为 $y_k$，集合 $S$ 对应的真实任务效用定义为
\begin{equation}
  u_k(S)=\ind\{S\ne\varnothing,\ c_k(S)=y_k\},
  \label{rq3:eq:true-utility}
\end{equation}
其中空集合效用为零，即未收到任何有效证据时不计作任务完成。需要说明的是，参考类别 $y_k$ 仅在离线训练与测试评价阶段可见，在线交付控制则采用第~\ref{rq3:subsec:utility} 节学习的效用网络进行估计。

在系统假设方面，考虑到单个观测任务的持续时间较短，假设任务内单跳无线接触保持连通、数据发送按串行方式执行、确认反馈占用独立控制信道，且各次发送的信道速率扰动相互独立。能耗统计涵盖数据帧与确认帧的无线发射支出，不计入候选元数据同步、端侧推理及无人机推进能耗。

底层链路沿用 LEADR 中的视距路径损耗、香农速率与发送能耗模型。由于数据串行传输且反馈信道独立，速率计算中不计同频并发数据干扰。结合节点到中心的距离 $d_{k,i}$、参考信道增益 $\rho_0$、有效带宽 $W_k$、噪声功率谱密度 $N_0$ 及数据发射功率 $P^{\mathrm d}$，候选 $m$ 所在链路的信道增益与预测速率分别为
\begin{equation}
  h_{k,i}=\frac{\rho_0}{\max\{d_{k,i},1\}^2},\qquad
  \bar r_{k,i}=W_k\log_2\left(1+\frac{P^{\mathrm d}h_{k,i}}{N_0W_k}\right).
  \label{rq3:eq:rate}
\end{equation}
上述预测速率在单个短任务内保持不变，距离与可用带宽则在任务间动态更新。

在物理层封装上，设平均每个 token 占 $b$ 字节，分组最大传输单元（Maximum Transmission Unit，MTU）的有效载荷上限为 $L_{\mathrm{MTU}}$ 字节，分组头部开销为 $H$ 字节，则候选 $m$ 实际需要传输的总比特数为
\begin{equation}
  B_m=8\left(b\ell_m+H\left\lceil\frac{b\ell_m}{L_{\mathrm{MTU}}}\right\rceil\right).
  \label{rq3:eq:bits}
\end{equation}
单次发送的预留时长记为
\begin{equation}
  \tau_{k,m}=\min\left\{\tau_{\max},\ \kappa\frac{B_m}{\bar r_{k,m}}\right\},
  \label{rq3:eq:duration}
\end{equation}
其中 $\kappa$ 为时间裕量系数，$\tau_{\max}$ 为单次发送时长上限，且整次传输必须落在任务剩余时间内。为刻画无线信道的随机波动，设第 $a$ 次尝试时的实际传输速率服从对数正态扰动：
\begin{equation}
  r_{k,m}^{(a)}=\bar r_{k,m}\exp\left(\sigma_r Z_a-\frac{\sigma_r^2}{2}\right),
  \qquad Z_a\sim\mathcal{N}(0,1),
  \label{rq3:eq:random-rate}
\end{equation}
其中 $\sigma_r$ 控制速率波动强度。若预留时长内能够传完全部 $B_m$ 比特，则该条证据完整送达；否则视为本次传输擦除，且依然耗尽完整预留时长。由此可得单次发送的成功概率：
\begin{equation}
  p_{k,m}=\Prb\{r_{k,m}^{(a)}\tau_{k,m}\ge B_m\}
  =\Phi\left(\frac{\ln(\bar r_{k,m}\tau_{k,m}/B_m)-\sigma_r^2/2}{\sigma_r}\right),
  \label{rq3:eq:success-probability}
\end{equation}
其中 $\Phi(\cdot)$ 为标准正态分布函数（$\sigma_r=0$ 时退化为确定性容量判定）。重传机制通过提供多次独立传输机会，提升证据至少成功送达一次的概率。

发送端依靠地面中心返回的累计确认（ACK）了解接收进展。如图~\ref{rq3:fig:delivery-scene} 所示，数据实际送达与发送端获知其送达往往并不同步：当确认尚在途中或已丢失时，继续重传可能造成不必要的重复占用，而盲目切换新候选又可能承担关键证据缺失的风险。

\begin{figure}[htbp]
\centering
% Illustrative events, without numerical timing or performance claims.
% Solid paths denote data; dashed paths denote cumulative feedback.
\begin{tikzpicture}[
  x=1mm,y=1mm,>=Stealth,
  font=\fontsize{9bp}{11.5bp}\selectfont,
  dline/.style={-{Stealth[length=1.6mm,width=1.2mm]},line width=.65pt,draw=blue!45!black},
  aline/.style={-{Stealth[length=1.6mm,width=1.2mm]},dashed,line width=.6pt,draw=orange!65!black},
  labelbg/.style={fill=white,inner sep=.4mm,align=center},
  pics/uav/.style={code={
    \draw[line width=.65pt] (-5.5,-2.5)--(5.5,2.5) (-5.5,2.5)--(5.5,-2.5);
    \foreach \x/\y in {-5.5/-2.5,-5.5/2.5,5.5/-2.5,5.5/2.5}{
      \draw[fill=black!3,line width=.45pt] (\x,\y) ellipse (3.4mm and 1.1mm);
      \draw[line width=.4pt] (\x-2.6,\y)--(\x+2.6,\y);
    }
    \draw[fill=black!12,line width=.5pt] (-2.2,-1.6) rectangle (2.2,1.6);
    \draw[fill=white,line width=.45pt] (-.9,-1.6) rectangle (.9,-3.5);
    \fill[black!65] (0,-2.8) circle (.45mm);
  }},
  pics/evidence/.style={code={
    \draw[fill=blue!4!white,line width=.4pt] (-4.5,-5.5) rectangle (4.5,5.5);
    \draw[line width=.35pt,draw=black!55] (-2.9,2.8)--(2.9,2.8)
      (-2.9,1.3)--(2.9,1.3) (-2.9,-.2)--(1.3,-.2);
  }}
]
\path[use as bounding box] (0,0) rectangle (148,-92);

% Two sensing UAVs and their candidate evidence.
\pic at (20,-16) {uav};
\pic at (68,-16) {uav};
\node at (20,-6) {感知无人机 A};
\node at (68,-6) {感知无人机 B};
\pic at (18,-32) {evidence};
\node at (18,-35) {$m_A$};
\pic at (65,-32) {evidence};
\node at (65,-35) {$m_B$};
\draw[draw=black!45,line width=.35pt] (19,-21)--(18,-26.5);
\draw[draw=black!45,line width=.35pt] (67,-21)--(65,-26.5);
\node at (18,-42) {已发送证据};
\node at (65,-42) {待发候选};

% Ground receiver: antenna and edge terminal, at natural scene proportions.
\draw[fill=black!2,line width=.45pt] (114,-49) rectangle (146,-52);
\draw[line width=.6pt] (124,-31)--(119,-49)--(129,-49)--cycle;
\draw[line width=.4pt] (122,-38)--(126,-38) (121,-43)--(127,-43);
\draw[line width=.4pt] (120,-47)--(126,-38) (128,-47)--(122,-38);
\fill[black!65] (124,-31) circle (.8mm);
\draw[line width=.45pt] (121.5,-28) arc[start angle=140,end angle=40,radius=3.2mm];
\draw[line width=.4pt] (119.8,-26.6) arc[start angle=140,end angle=40,radius=5.4mm];
\draw[fill=black!5,line width=.5pt] (132,-40) rectangle (143,-46.5);
\draw[line width=.5pt] (137.5,-46.5)--(137.5,-49);
\draw[line width=.5pt] (134,-49)--(141,-49);
\node at (129.5,-56) {地面接收中心};
\draw[draw=black!35,line width=.4pt] (3,-53)--(146,-53);

% An already delivered message and a prospective alternative.
\draw[dline] (29,-18.5) to[out=-7,in=166] (119,-34);
\node[labelbg] at (84,-28) {$m_A$ 已到达};
\draw[dline,draw=green!40!black] (77,-18) to[out=0,in=155] (118,-27);
\node[labelbg] at (101,-16) {可新发 $m_B$};
\draw[aline] (116,-46.5) -- (36,-46.5) to[out=180,in=-40] (28,-23);
\node[labelbg] at (80,-48) {累计确认延迟或丢失};

% A short event view separates sender knowledge from receiver reality.
\draw[draw=black!25,line width=.4pt] (0,-60)--(148,-60);
\node[anchor=east] at (24,-68) {发送端};
\node[anchor=east] at (24,-81) {接收端};
\draw[->,draw=black!55,line width=.4pt] (27,-68)--(146,-68);
\draw[->,draw=black!55,line width=.4pt] (27,-81)--(146,-81);
\node[anchor=north east] at (146,-85) {时间};
\draw[dline] (36,-68)--(57,-81);
\fill[blue!45!black] (36,-68) circle (.6mm);
\fill[blue!45!black] (57,-81) circle (.6mm);
\node[labelbg] at (45,-63.5) {发送 $m_A$};
\node[labelbg] at (52,-87) {$m_A$ 已接收};
\draw[aline] (57,-81)--(86,-69.5);
\draw[draw=orange!65!black,line width=.7pt] (84.8,-70.7)--(87.2,-68.3)
  (84.8,-68.3)--(87.2,-70.7);
\node[labelbg] at (74,-77) {确认未到};
\draw[draw=black!50,line width=.4pt] (96,-62)--(96,-88.5);
\node[labelbg] at (96,-63.5) {当前};
\node[labelbg] at (121,-73) {补发 $m_A$ 或新发 $m_B$};
\node[labelbg] at (120,-86) {实际已收 $\{m_A\}$\\重复不新增证据};
\end{tikzpicture}

\bicaption{语义数据送达与累计确认不同步的通信示意}{Semantic data delivery with delayed or missing cumulative acknowledgment}
\label{rq3:fig:delivery-scene}
\end{figure}

具体而言，每次数据发送结束后，接收中心均会生成一条包含尝试序号 $j$ 与当前接收集合快照 $S_j^{\mathrm{ack}}$ 的累计确认帧（短控制头可独立标识调度序号，因此即便本次数据载荷解码失败，中心仍可回传最新接收位图，以便发送端同步确认已收集合与漏收情况）。记第 $j$ 次数据发送的结束时刻为 $t_j^{\mathrm{end}}$，反馈帧长度为 $B^{\mathrm f}$ 比特，则各确认帧沿对应链路按预测速率串行发出的完成时刻满足
\begin{equation}
  v_j=\max\{t_j^{\mathrm{end}},v_{j-1}\}+\tau_j^{\mathrm f},
  \qquad \tau_j^{\mathrm f}=\frac{B^{\mathrm f}}{\bar r_{k,m_j}},
  \label{rq3:eq:feedback-serialization}
\end{equation}
其中初始值 $v_{-1}=0$。串行发出后，每条确认帧以概率 $\epsilon_{\mathrm f}$ 发生丢失；若未丢失，则经历随机额外时延 $L_j$ 到达发送端：
\begin{equation}
  L_j\sim\mathcal{U}\left(\frac{\bar d_{\mathrm f}}{2},\frac{3\bar d_{\mathrm f}}{2}\right),
  \qquad t_j^{\mathrm{ack}}=v_j+L_j.
  \label{rq3:eq:feedback-delay}
\end{equation}
数据成败、反馈丢失及反馈延迟彼此独立抽样。由于确认帧可能乱序到达，控制器仅保留尝试序号最大的有效快照，不会用陈旧位图覆盖较新状态。

记时刻 $t$ 发送端掌握的全部可见历史为 $\history_{k,t}$，涵盖候选与链路预测、过往发送记录、已到达的累计确认以及当前剩余时间与电量（不包含真实的隐藏接收集合和在途确认）。定义接收信念为当前真实接收集合的条件概率分布：
\begin{equation}
  b_{k,t}(S)=\Prb\{S_{k,t}=S\mid\history_{k,t}\},\qquad S\in\states_k.
  \label{rq3:eq:belief-definition}
\end{equation}
在数据成败与反馈丢失相互独立的前提下，即便尚未收到确认，发送端仍可结合历史发送记录推断各接收状态的概率，而主动等待能否换取新状态信息则取决于在途确认的到达概率。

\subsection{通信资源与优化目标}
\label{rq3:subsec:energy}
\label{rq3:subsec:objective}

设感知节点 $i$ 的初始通信电量为 $E_i^{(0)}$，则单次数据发送与对应确认回传的发射能耗分别为
\begin{equation}
  e_{k,m}^{\mathrm d}=P^{\mathrm d}\tau_{k,m},\qquad
  e_{k,m}^{\mathrm f}=P^{\mathrm f}\tau_{k,m}^{\mathrm f}.
  \label{rq3:eq:transmit-energy}
\end{equation}
其中数据发送能耗从对应感知节点 $s(m)$ 电量中逐次扣减；接收中心独立供电，但反馈帧一经生成（无论后续是否丢失或超时到达），其发射能耗 $e_{k,m}^{\mathrm f}$ 均计入系统总通信成本。记节点 $i$ 在第 $k$ 项任务中累计消耗的数据发送能量为 $e_i(k)$，则跨任务尺度的节点剩余电量演化为
\begin{equation}
  E_i(k+1)=E_i(k)-e_i(k),\qquad E_i(1)=E_i^{(0)}.
  \label{rq3:eq:energy-evolution}
\end{equation}

在每个决策时刻，控制器可执行三类动作：发送某条候选 $m$（若该候选此前已发过即为重传）、主动等待时长 $w$ 或直接停止当前任务。记当前已执行尝试次数为 $n_{k,t}$，单任务尝试上限为 $A_{\max}$，则同时满足剩余时限、节点电量与尝试上限的可行发送集合为
\begin{equation}
  \begin{split}
  \actions^{\mathrm{send}}_{k,t}=\{m\in\mathcal{M}_k:\;&
  \tau_{k,m}\le D_k-t,\\[-2pt]
  &e_{k,m}^{\mathrm d}\le E_{s(m)}(k,t),\quad n_{k,t}<A_{\max}\}.
  \end{split}
  \label{rq3:eq:action-mask}
\end{equation}
等待动作须满足 $0<w\le D_k-t$，停止动作则始终可行。通过式~\eqref{rq3:eq:action-mask} 的硬约束校验，系统严格杜绝超时发送或电量透支。

记策略 $\pi$ 在任务 $k$ 上的停止时刻为 $T_k$、终止时的真实接收集合为 $S_k^{\mathrm{end}}$、实际执行的数据发送尝试集合为 $\mathcal{T}_k$，则该任务的总通信能耗为
\begin{equation}
  E_k^{\mathrm{comm}}=\sum_{a\in\mathcal{T}_k}
    \left(e_{k,m_a}^{\mathrm d}+e_{k,m_a}^{\mathrm f}\right).
  \label{rq3:eq:total-communication-energy}
\end{equation}
在有限任务序列中，交付控制的目标是在提升按时正确完成率与压缩通信能耗之间取得平衡，相应的全局随机优化问题可表述为
\begin{equation}
  \begin{aligned}
    \max_{\pi}\quad &\frac{1}{N_{\mathrm{task}}}\E_{\pi}\left[
      \sum_{k=1}^{N_{\mathrm{task}}}\left(u_k(S_k^{\mathrm{end}})-\lambda_E
      \frac{E_k^{\mathrm{comm}}}{E_{\mathrm{ref}}}\right)\right] \\
    \text{s.t.}\quad &\sum_{k=1}^{N_{\mathrm{task}}} e_i(k)\le E_i^{(0)},\quad i\in\mathcal{N},\\
    &0\le T_k\le D_k,\quad |\mathcal{T}_k|\le A_{\max},\\
    &a_{k,t}\text{ 仅依赖 }\history_{k,t}\text{，且满足式~\eqref{rq3:eq:action-mask}},
  \end{aligned}
  \label{rq3:eq:global-problem}
\end{equation}
其中 $E_{\mathrm{ref}}>0$ 为能量归一化尺度，$\lambda_E\ge0$ 为能耗惩罚系数。该问题的难点在于接收状态隐藏、真实任务效用不可观测以及多任务共享有限电量三者交织：多发一次数据既可能补齐当前判断证据，也会直接压缩后续任务的能量余量。为此，下节先利用离线学习的效用网络替代不可观测的真实效用，再基于接收信念、跨任务虚拟能量队列与有限前瞻，构造式~\eqref{rq3:eq:global-problem} 的在线近似求解方法。本章主要符号汇总于表~\ref{rq3:tab:notation}。

\begin{table}[t]
  \centering\small
  \bicaption{语义交付控制模型的主要符号}{Main notation of the semantic delivery control model}
  \label{rq3:tab:notation}
  \begin{tabularx}{\linewidth}{@{}lY@{}}
    \toprule
    符号 & 含义\\
    \midrule
    $N,\ N_{\mathrm{task}}$ & 感知节点数及连续执行的任务总数\\
    $\mathcal{M}_k,\ K_k$ & 任务候选证据集合及其规模\\
    $S_{k,t},\ b_{k,t}(S)$ & 实际接收集合及其条件概率\\
    $c_k(S),\ u_k(S)$ & 接收集合对应的预测类别与真实任务效用\\
    $\hat q_\theta(S),\ q_k(S)$ & 原始预测效用与一致性修正效用\\
    $\bar r_{k,m},\ p_{k,m}$ & 预测链路速率与单次送达概率\\
    $\tau_{k,m},\ D_k$ & 数据发送时长与任务截止时间\\
    $\bar d_{\mathrm f},\epsilon_{\mathrm f}$ & 反馈平均额外时延与反馈丢失概率\\
    $E_i(k,t),\ Z_i(k)$ & 节点剩余发送能量与虚拟能量队列\\
    $\pi_{k,m}$ & 一次发送及其反馈的归一化资源价格\\
    $\eta_E,\eta_T$ & 局部控制的能量权重与时间权重\\
    \bottomrule
  \end{tabularx}
\end{table}
\FloatBarrier
\section{一致性感知的信念反馈控制方法}
\label{rq3:sec:method}

如图~\ref{rq3:fig:controller} 所示，C-BFC 的在线运行划分为任务初始化与任务内滚动控制两个阶段。在任务开始时，控制器一次性枚举所有可能的非空接收集合并估计其任务效用，再经同类别一致性投影生成静态效用表；进入传输阶段后，控制器只需随事件推进更新接收信念、剩余时限与能量状态，在线权衡发送、等待与停止动作。整个重传周期内既无需重算效用表，也不必重新调用视觉语言模型，从而保持极低的在线计算开销。

\begin{figure}[t]
  \centering
  % The task utility table is built once; only states and action values recur.
\begin{tikzpicture}[x=1mm,y=1mm,
  font=\fontsize{9bp}{11.5bp}\selectfont,
  cbbox/.style={draw,line width=.45pt,fill=white,align=center,inner sep=.8mm},
  cbarrow/.style={-{Stealth[length=1.6mm,width=1.2mm]},line width=.5pt},
  cbgroup/.style={draw,line width=.45pt},
  cbtitle/.style={font=\fontsize{9bp}{11.5bp}\selectfont\bfseries}]
  \path[use as bounding box] (0,0) rectangle (148,-93);
  \draw[cbgroup] (0,0) rectangle (48,-48);
  \draw[cbgroup] (52,0) rectangle (145,-48);
  \node[cbtitle] at (24,-3.2) {任务效用表构造};
  \node[cbtitle] at (98.5,-3.2) {在线状态更新};

  \node[cbbox,text width=42mm,minimum height=11mm] (sets) at (24,-13)
    {候选证据 $\mathcal M_k$，$K\le4$\\枚举 $2^K$ 个接收集合};
  \node[cbbox,text width=42mm,minimum height=11mm] (network) at (24,-28)
    {效用网络（非空集合）\\$29\rightarrow128\rightarrow96\rightarrow1$};
  \node[cbbox,text width=42mm,minimum height=11mm] (utility) at (24,-42)
    {分类一致性投影\\$q_k(S)$，$q_k(\varnothing)=0$};
  \draw[cbarrow] (sets.south) -- (network.north);
  \draw[cbarrow] (network.south) -- (utility.north);

  \node[cbbox,text width=42mm,minimum height=18mm] (history) at (74.5,-17)
    {累计确认\\已发送记录\\送达率 $p_{k,m}$};
  \node[cbbox,text width=42mm,minimum height=18mm] (resources) at (122.5,-17)
    {剩余时限\\节点电量与虚拟队列\\$h,E_i,Z_i$};
  \node[cbbox,text width=42mm,minimum height=17mm] (belief) at (74.5,-39)
    {累计快照及后续记录\\接收状态滤波\\$b_{k,t}(S)$};
  \node[cbbox,text width=42mm,minimum height=17mm] (price) at (122.5,-39)
    {时限、电量与尝试上限\\动作可行性\\资源价格 $\pi_{k,m}$};
  \draw[cbarrow] (history.south) -- (belief.north);
  \draw[cbarrow] (resources.south) -- (price.north);

  \node[cbbox,text width=140mm,minimum height=7.5mm,fill=black!3] (expected) at (74,-54)
    {期望效用 $U_k(b)=\sum_S b(S)q_k(S)$；结合资源价格评价可行动作};
  \foreach \source in {utility,belief,price}
    \draw[cbarrow] (\source.south) -- (expected.north -| \source.south);

  \node[cbbox,text width=42mm,minimum height=13mm] (send) at (24,-68.5)
    {两步发送前瞻：新发/重传\\$V_m^{\mathrm{send}}(b;h)$};
  \node[cbbox,text width=42mm,minimum height=13mm] (wait) at (74.5,-68.5)
    {反馈等待：至多三个时长\\$V^{\mathrm{wait}}(w;b,h)$};
  \node[cbbox,text width=42mm,minimum height=13mm] (stop) at (125,-68.5)
    {停止当前任务\\净增益为零};
  \draw[line width=.5pt] (expected.south) -- (74,-59.8);
  \draw[line width=.5pt] (24,-59.8) -- (125,-59.8);
  \foreach \x/\target in {24/send,74.5/wait,125/stop}
    \draw[cbarrow] (\x,-59.8) -- (\target.north);

  \node[cbbox,text width=138mm,minimum height=12mm,fill=black!3] (action) at (74,-85)
    {执行最大正收益的发送/等待动作，否则停止\\
     未停止则更新历史、时间及电量；任务结束后更新 $Z_i(k+1)$};
  \foreach \x/\source in {24/send,74.5/wait,125/stop}
    \draw[line width=.5pt] (\source.south) -- (\x,-77);
  \draw[line width=.5pt] (24,-77) -- (125,-77);
  \draw[cbarrow] (74,-77) -- (action.north);
  % Dashed return enters the online-state group, not the fixed utility table.
  \draw[cbarrow,dashed] (action.east) -- (147.5,-85) -- (147.5,-3.2) -- (145,-3.2);
\end{tikzpicture}

  \bicaption{C-BFC 的整体计算流程}{Overall computation flow of C-BFC}
  \label{rq3:fig:controller}
\end{figure}

\subsection{消息组合的任务收益估计}
\label{rq3:subsec:utility}

为了在不依赖真实参考标签的情况下评估任意非空接收集合 $S$ 的判别质量，本文构造了 29 维集合特征向量 $\bm f_k(S)\in\mathbb R^{29}$，涵盖五类融合后语义信念、最大置信度、前二类别间隔、分布熵、证据条数与总长度、发送方置信度、局部观测质量、节点间一致程度以及各表达粒度占比。该特征仅以集合 $S$ 已经成功送达为条件描述语义内容质量，无线链路丢包风险则解耦交由后续接收信念滤波处理。

利用训练集特征的均值 $\bm\mu$ 与标准差 $\bm s$，对输入特征按维执行标准化：
\begin{equation}
  \widetilde{\bm f}_k(S)=\frac{\bm f_k(S)-\bm\mu}{\max\{\bm s,0.1\}}.
  \label{rq3:eq:feature-normalization}
\end{equation}
效用估计器采用轻量级两层全连接网络，结构为
\begin{equation}
  29\longrightarrow128\longrightarrow96\longrightarrow1,
  \label{rq3:eq:network-architecture}
\end{equation}
共包含 16,321 个可训练参数。隐藏层选用 Sigmoid 线性单元（Sigmoid Linear Unit，SiLU）激活函数，训练阶段施加 0.1 概率的随机失活（Dropout），输出端经 Sigmoid 映射得到原始效用估计 $\hat q_\theta(S)\in(0,1)$。

离线训练时，枚举训练任务中的所有非空接收集合，以式~\eqref{rq3:eq:true-utility} 的真实分类正确性为监督标签，最小化二元交叉熵损失：
\begin{equation}
  \mathcal L(\theta)=-\frac{1}{n_{\mathcal D}}
  \sum_{(k,S)\in\mathcal D}
  \left[u_k(S)\ln\hat q_\theta(S)
  +(1-u_k(S))\ln\bigl(1-\hat q_\theta(S)\bigr)\right],
  \label{rq3:eq:utility-loss}
\end{equation}
其中 $n_{\mathcal D}$ 为样本总数。训练完毕后选取验证集交叉熵最低的模型参数并冻结，在线推理阶段不再需要参考标签。

然而，神经网络对各集合独立打分时存在一个隐蔽缺陷：同一任务中，哪怕两个非空集合 $S$ 与 $S'$ 融合后输出完全相同的预测类别（即 $c_k(S)=c_k(S')$），神经网络仍可能因置信度高低不同给出 $\hat q_\theta(S')>\hat q_\theta(S)$ 的估值。按照任务完成定义，只要预测类别一致，二者的真实正确性必然完全相同（$u_k(S)=u_k(S')$）；这种虚假的估值差会误导控制器为了追求“更高的信心”而白白重发或补发根本无法翻转最终判断的冗余消息。

为此，本文提出同类别一致性投影。对任务 $k$ 中所有输出同一预测类别 $c$ 的非空集合划分等价组：
\begin{equation}
  \mathcal G_{k,c}=\{S\in\states_k:S\ne\varnothing,\ c_k(S)=c\}.
  \label{rq3:eq:class-groups}
\end{equation}
在最小化与原始估值平方偏差的前提下，强制同组集合取相等效用：
\begin{equation}
  \begin{aligned}
    \min_{\{q_k(S)\}}\quad &\sum_{S\ne\varnothing}
       \bigl(q_k(S)-\hat q_\theta(S)\bigr)^2\\
    \mathrm{s.t.}\quad &q_k(S)=q_k(S'),\quad S,S'\in\mathcal G_{k,c},
  \end{aligned}
  \label{rq3:eq:consistency-projection}
\end{equation}
空集合则单独固定为 $q_k(\varnothing)=0$。该凸优化问题不仅具有简洁的闭式解，还能严格保证不增大估计的平方误差。

\begin{proposition}[一致性投影的闭式解与误差性质]
\label{rq3:prop:projection}
对于任一非空类别组，式~\eqref{rq3:eq:consistency-projection} 的唯一最优解为组内原始估计效用的算术平均：
\begin{equation}
  q_k(S)=\frac{1}{|\mathcal G_{k,c_k(S)}|}
    \sum_{S'\in\mathcal G_{k,c_k(S)}}\hat q_\theta(S').
  \label{rq3:eq:projection-solution}
\end{equation}
记 $\bm u_k$ 为各非空集合的真实效用向量，则投影后的效用向量 $\bm q_k$ 满足
\begin{equation}
  \|\bm q_k-\bm u_k\|_2^2
  \le\|\widehat{\bm q}_k-\bm u_k\|_2^2.
  \label{rq3:eq:projection-error}
\end{equation}
\end{proposition}

\begin{proof}
对任意类别组 $\mathcal G_{k,c}$，一致性约束等价于令组内所有分量共享同一标量 $z_c$，对应子目标 $\sum_{S\in\mathcal G_{k,c}}(z_c-\hat q_\theta(S))^2$ 关于 $z_c$ 严格凸，令其一阶导数为零即得唯一最优解~\eqref{rq3:eq:projection-solution}。进一步注意到同组集合具有相同的真实效用常数 $u_c$，利用均值中心化展开可得
\begin{equation}
  \sum_{S\in\mathcal G_{k,c}}(\hat q_\theta(S)-u_c)^2
  =\sum_{S\in\mathcal G_{k,c}}(\hat q_\theta(S)-z_c)^2
   +|\mathcal G_{k,c}|(z_c-u_c)^2.
  \label{rq3:eq:projection-decomposition}
\end{equation}
由于右侧第一项恒非负，对所有类别组求和即证得式~\eqref{rq3:eq:projection-error}。
\end{proof}

需要指出的是，上述投影仅依据接收中心自身可计算的预测类别 $c_k(S)$ 分组，在线阶段完全无需真实标签，且修正后的 $q_k(S)$ 表示各集合的任务效用得分（无需跨类别归一化）。同时，即便两个集合因预测相同而获得一致效用，各自对应的传输路径和送达概率依然存在差异，这部分差异会通过接收信念 $b_{k,t}$ 完整传递到后续动作决策中。效用模型训练与任务效用表构造流程见算法~\ref{rq3:alg:utility}。

\begin{algorithm}[H]
  \bicaption{语义效用模型训练与任务效用表构造}{Semantic utility learning and task utility table construction}
  \label{rq3:alg:utility}
  \small
  \begin{algorithmic}[1]
    \Require 训练与验证语义任务；固定候选生成与融合规则
    \Ensure 效用网络参数 $\theta$；任务 $k$ 的效用表 $q_k$
    \For{每个训练任务及其非空接收集合 $S$}
      \State 提取 $\bm f_k(S)$，按式~\eqref{rq3:eq:true-utility} 计算监督目标
    \EndFor
    \State 计算训练特征统计量，按式~\eqref{rq3:eq:utility-loss} 优化网络
    \State 选择验证交叉熵最低的网络参数并冻结
    \For{待执行任务的每个非空集合 $S$}
      \State 推理 $\hat q_\theta(S)$，计算接收端预测类别 $c_k(S)$
    \EndFor
    \State 按预测类别分组，利用式~\eqref{rq3:eq:projection-solution} 计算 $q_k(S)$
    \State 令 $q_k(\varnothing)=0$
    \State \Return $\theta$ 与 $q_k$
  \end{algorithmic}
\end{algorithm}

\subsection{接收状态估计}
\label{rq3:subsec:filter}

在不考虑新确认到达的情况下，发送一次候选 $m$ 会以概率 $p_{k,m}$ 使接收集合由 $S$ 扩充为 $S\cup\{m\}$，以概率 $1-p_{k,m}$ 保持不变。由此定义信念前向转移算子：
\begin{equation}
  (\mathcal T_m b)(S')=(1-p_{k,m})b(S')
   +p_{k,m}\sum_{S:S\cup\{m\}=S'}b(S).
  \label{rq3:eq:belief-transition}
\end{equation}
若原集合 $S$ 已包含 $m$，则无论本次发送成功与否集合均不改变，从而自然刻画了“已收证据重复发送无增量信息”的性质。

每当控制器进行决策时，先检查迄今已收到的最新累计确认。若最新有效确认携带了第 $j$ 次尝试结束时的快照 $S_j^{\mathrm{ack}}$，则先将第 $j$ 次尝试后的状态锚定为单点分布 $\delta_{S_j^{\mathrm{ack}}}$，再依次重放其后截至第 $n$ 次的所有未确认发送记录：
\begin{equation}
  b_{k,t}=\mathcal T_{m_n}\cdots\mathcal T_{m_{j+1}}
          \delta_{S_j^{\mathrm{ack}}}.
  \label{rq3:eq:ack-replay}
\end{equation}
若尚未收到任何确认，则直接从初始空集合 $\delta_{\varnothing}$ 重放全部发送记录；乱序迟到的陈旧确认一律忽略。

\begin{proposition}[接收信念的有效性]
\label{rq3:prop:belief}
对于任意非负且归一化的信念分布 $b$，经式~\eqref{rq3:eq:belief-transition} 转移后的结果仍构成合法的概率分布。在第~\ref{rq3:subsec:channel} 节的数据传输与反馈独立性假设下，式~\eqref{rq3:eq:ack-replay} 精确给出了给定最新累计确认与后续发送历史时当前接收集合的后验条件分布。
\end{proposition}

\begin{proof}
式~\eqref{rq3:eq:belief-transition} 中各项非负，且每个原状态 $S$ 的概率质量以权重 $1-p_{k,m}$ 与 $p_{k,m}$ 分别流向 $S$ 与 $S\cup\{m\}$，总概率质量守恒为 1。累计确认精确锁定了第 $j$ 次尝试结束时的真实集合；其后各次数据发送相互条件独立，且反馈的延迟或丢失独立于数据解码结果（即“确认尚未到达”本身不泄露数据失败的额外信息），因此顺次施加转移算子即得当前精确的条件分布。
\end{proof}

直观来看，若某条单次成功概率为 $p$ 的证据已连续独立发送 $r$ 次，即使尚未收到确认，其至少送达一次的概率也已升至 $1-(1-p)^r$。借助接收信念滤波，控制器能够自然捕捉到连续重传的边际收益递减，避免了传统方法把“未收到确认”一概等同于“接收失败”所引发的过度重传。结合一致性效用表 $q_k$，当前接收信念对应的期望任务效用为
\begin{equation}
  U_k(b)=\sum_{S\in\states_k}b(S)q_k(S).
  \label{rq3:eq:expected-utility}
\end{equation}

\subsection{能量约束下的发送与等待决策}
\label{rq3:subsec:queue}
\label{rq3:subsec:decision}
\label{rq3:subsec:complexity}

为防止前序任务过度耗电导致后续任务无电可用，沿用第~\ref{rq1:chapter} 章的李雅普诺夫虚拟队列思想，并将时间尺度由网络时隙切换为连续任务序列。每个感知节点维护跨任务虚拟能量队列 $Z_i(k)$，记录相对平均任务额度的累积超支压力：
\begin{equation}
  Z_i(k+1)=\max\left\{0,Z_i(k)+e_i(k)-\frac{E_i^{(0)}}{N_{\mathrm{task}}}\right\},
  \qquad Z_i(1)=0.
  \label{rq3:eq:virtual-queue}
\end{equation}
当某节点在前期任务中频繁超额发送时，$Z_i(k)$ 持续累积升高，而在低消耗任务后逐步回落。同时，结合集群当前的剩余总电量与待完成任务数，动态计算后续每项任务的平均可用能量额度：
\begin{equation}
  A_{k,t}=\frac{\sum_{i\in\mathcal{N}}E_i(k,t)}{N_{\mathrm{task}}-k+1}.
  \label{rq3:eq:remaining-allowance}
\end{equation}
由此定义一次候选发送（含对应确认反馈）的归一化资源价格：
\begin{equation}
  \pi_{k,m}=\frac{
      \left(1+\dfrac{Z_{s(m)}(k)}{E_{s(m)}^{(0)}}\right)e_{k,m}^{\mathrm d}
      +e_{k,m}^{\mathrm f}}
      {\max\{A_{k,t},\varepsilon_E\}},
  \label{rq3:eq:resource-price}
\end{equation}
其中小常数 $\varepsilon_E>0$ 用于避免数值除零。分子同时计入数据与反馈能耗，并对历史透支较重的节点施加额外溢价（反馈能耗由地面中心承担，不扣减无人机电量）；分母则反映剩余任务的平均能量余量。在实际控制中，硬电量约束仍由动作掩码~\eqref{rq3:eq:action-mask} 严格把关，而虚拟队列通过动态抬升 $\pi_{k,m}$ 引导前期耗电偏快的节点在后续任务中趋于克制。

为简化表达，下文省去任务与时刻下标，记任务剩余时间为 $h=D_k-t$。在当前信念 $b$ 下，发送候选 $m$ 带来的单步即时净收益定义为
\begin{equation}
  g_m(b;h)=U_k(\mathcal T_m b)-U_k(b)
     -\eta_E\pi_{k,m}-\eta_T\frac{\tau_{k,m}}{D_k},
  \label{rq3:eq:send-gain}
\end{equation}
其中 $\eta_E$ 与 $\eta_T$ 分别为能量价格权重与显式时间代价权重；凡违反时限或节点电量约束的候选，其收益直接置为 $-\infty$。值得注意的是，式~\eqref{rq3:eq:send-gain} 统一评估新发与重传动作：重传通过拉高已发证据的送达概率获取收益，新发则通过引入互补证据拓展集合组合，两者在统一框架下直接比拼净收益，无需人为设定固定最大重传次数。

由于多条互补证据常常存在“单条发送无法翻转类别、两条组合后才产生显著增益”的协同效应，仅看单步即时收益容易过早停止。为此，控制器向前展开两步发送前瞻：
\begin{equation}
  V_m^{\mathrm{send}}(b;h)=g_m(b;h)
    +\max\left\{0,\max_{m'\in\mathcal F(m)}
      g_{m'}(\mathcal T_m b;h-\tau_{k,m})\right\},
  \label{rq3:eq:two-step-lookahead}
\end{equation}
其中第二步可行集 $\mathcal F(m)$ 依据执行 $m$ 后的剩余时间、剩余电量与尝试上限构造（前瞻计算内固定当前能量价格，实际反馈留待下次决策再行消化；若多个候选两步收益打平，则依次按单步即时收益与固定候选序号打破平局）。

与立即发送相对，主动等待虽然不发出新数据包，却有机会在等待窗口内等到在途累计确认，从而消除接收状态的不确定性、避免盲目重传。为控制分支规模，本文围绕最新一条尚未返回的确认快照评估等待价值。设该确认帧的串行发送结束时刻为 $v_j$，截至当前的在途年龄为 $a=t-v_j$，额外时延累积分布函数为 $F(\cdot)$，则在已知当前尚未收到该确认的条件下，继续等待时长 $w$ 内收到它的条件概率为
\begin{equation}
  \rho(w)=\frac{(1-\epsilon_{\mathrm f})[F(a+w)-F(a)]}
                {1-(1-\epsilon_{\mathrm f})F(a)},
  \label{rq3:eq:conditional-feedback}
\end{equation}
其中分母同时包含了“尚在途中”和“已经丢失”两种未达情形；一旦 $a$ 超出时延分布上界使分子为零，继续等待该快照的信息价值便自动归零。

记再执行至多一次发送所能获得的最大净收益为
\begin{equation}
  G(b;h)=\max\left\{0,\max_m g_m(b;h)\right\}.
  \label{rq3:eq:one-step-value}
\end{equation}
若等待 $w$ 后成功收到确认，接收集合将坍缩为某个确切状态 $S$（分布变为点质量 $\delta_S$）；若未收到，则沿用原信念 $b$。据此，等待时长 $w$ 的近似信息价值可写为
\begin{equation}
  \begin{split}
    V^{\mathrm{wait}}(w;b,h)
      ={}&\rho(w)\sum_{S\in\states_k}b(S)G(\delta_S;h-w)\\
         &+[1-\rho(w)]G(b;h-w)-\eta_T\frac{w}{D_k}.
  \end{split}
  \label{rq3:eq:waiting-value}
\end{equation}
需要强调的是，即便显式时间权重 $\eta_T=0$，等待消耗的时间依然通过缩短剩余窗口 $h-w$ 压缩了后续可行发送空间，从而天然体现出等待的机会成本。具体执行时，等待候选集仅需考察三个具有明确物理意义的时长节点——固定短等待 $w_0$、等待至反馈平均额外时延以及等待至时延分布上界：
\begin{equation}
  \mathcal W_t=\left\{w_0,\ \max(w_0,\bar d_{\mathrm f}-a),\
     \max(w_0,1.5\bar d_{\mathrm f}-a)\right\}\cap(0,h).
  \label{rq3:eq:wait-candidates}
\end{equation}
在每次决策中，控制器比较所有可行发送动作的 $V_m^{\mathrm{send}}$ 与等待候选的 $V^{\mathrm{wait}}$，执行正收益最大的动作；若所有动作收益均不大于零，则果断停止当前任务，并在任务结束后更新虚拟能量队列。完整在线流程见算法~\ref{rq3:alg:cbfc}。

\begin{algorithm}[t]
  \bicaption{C-BFC 在线语义交付控制}{Online semantic delivery control of C-BFC}
  \label{rq3:alg:cbfc}
  \small
  \begin{algorithmic}[1]
    \Require 候选与链路预测；效用模型 $\theta$；初始节点能量；$\eta_E,\eta_T$
    \Ensure 每个任务的发送、等待与停止动作序列
    \State 初始化 $E_i(1)=E_i^{(0)}$，$Z_i(1)=0$
    \For{任务 $k=1,\ldots,N_{\mathrm{task}}$}
      \State 根据算法~\ref{rq3:alg:utility} 构造 $q_k$，令 $t=0$，$b=\delta_{\varnothing}$
      \While{任务尚未停止且 $t<D_k$}
        \State 读取已到达的累计确认与可见发送历史
        \State 按式~\eqref{rq3:eq:ack-replay} 更新接收信念 $b$
        \State 构造可行动作，并按式~\eqref{rq3:eq:resource-price} 计算能量价格
        \If{不存在可行发送动作}
          \State 停止当前任务
        \Else
          \State 按式~\eqref{rq3:eq:two-step-lookahead} 评价各发送动作
          \State 按式~\eqref{rq3:eq:conditional-feedback}--\eqref{rq3:eq:waiting-value} 评价等待候选
          \If{发送与等待的最大收益不大于零}
            \State 停止当前任务
          \Else
            \State 执行收益最大的发送或等待动作
            \State 更新当前时间、发送能量与可见历史
          \EndIf
        \EndIf
      \EndWhile
      \State 汇总节点发送支出，按式~\eqref{rq3:eq:virtual-queue} 更新 $Z_i(k+1)$
    \EndFor
  \end{algorithmic}
\end{algorithm}

就计算复杂度而言，设单任务候选数为 $K$，效用网络参数量为 $P_\theta$。任务初始化阶段批量推理至多 $2^K$ 个集合的效用；在线控制中，单次信念转移与期望效用求值复杂度为 $O(2^K)$，两步发送前瞻为 $O(K^2 2^K)$，而等待评价需枚举反馈可能揭示的接收集合及随后的单步最优发送，因此单次决策的总体时间复杂度为
\begin{equation}
  O\left(K^2 2^K+|\mathcal W_t|K4^K\right).
  \label{rq3:eq:complexity}
\end{equation}
得益于前端 EviComRL 的预筛选，本章中 $K\le4$（状态空间至多 16 个集合，等待候选至多 3 个），可直接精确枚举，单步决策耗时极低。若未来应用场景中候选规模大幅扩张导致状态数呈指数增长，则需进一步引入候选剪枝或因子化近似滤波；此外，本文的同类别一致性投影立足于离散分类任务，面向开放式文本生成或连续值回归任务时，还需结合相应的任务损失函数重新构造一致性约束。

\FloatBarrier
\section{实验设置与评价方法}
\label{rq3:sec:experiments}

\subsection{数据与实现参数}
\label{rq3:subsec:experiment-data}

实验采用由 VQAv2/COCO 图像构造的预先生成的观测数据\citep{coco,vqa-v2}。每幅图像及查询构成一个任务，局部裁剪提供九个节点的观测，中心不访问原图。离线视觉语言模型生成局部证据与全图参考标签，分为清晰可见、大部分可见、不确定、大部分遮挡和完全遮挡五类。以接收端分类与全图参考标签一致作为正确判据，比较给定语义内容在不同交付策略下的任务表现。

候选由训练完成的 EviComRL 在给定 token 预算下生成。该模块用策略网络选择候选消息，并用价值网络评价后续收益，采用 PPO 训练\citep{ppo}。生成后按固定次序合并去重，保留至多四条。实验中固定该策略的参数，所有交付方法使用相同的候选内容和次序，以比较相同语义内容下的交付控制效果。

训练、验证和测试分区分别含 600、595、611 个任务，图像互不重叠。枚举训练任务的非空集合生成监督样本，同一图像的集合不跨分区。效用网络采用带解耦权重衰减的自适应矩估计优化器 AdamW，学习率 $8\times10^{-4}$、权重衰减 0.02、批量 256，最多训练 60 轮，以验证交叉熵选择参数。

控制参数在参考条件、12 ms 反馈延迟、60\% 反馈丢失及 12 ms 任务时限四种验证条件下选择：优先平均正确完成率，并列时选择低能耗。所得参数用于全部测试。C-BFC 三个训练种子均取 $\eta_E=0.24$、$\eta_T=0$；时限仍通过可行域和等待机会成本作用于决策。

学习方法使用三个训练种子与三个网络种子的九种组合，固定协议使用三个网络种子。每次按相同次序处理 611 个任务，电量连续扣减；各方法对同一任务、候选和尝试使用相同随机条件。误差线为重复实验标准差，重复实验改变初始化及网络条件，视觉样本固定。

\label{rq3:subsec:experiment-settings}

参数见表~\ref{rq3:tab:parameters}。节点距离在 100--260 m 内平滑变化，可用带宽系数为 0.25--1.0；按式~\eqref{rq3:eq:rate} 预测速率，按式~\eqref{rq3:eq:success-probability} 抽样送达。每任务最多九次发送，来自同一节点的不同候选共用电量。

\begin{table}[t]
  \centering\small
  \bicaption{语义交付实验的主要通信与控制参数}{Main communication and control parameters for semantic delivery}
  \label{rq3:tab:parameters}
  \begin{tabularx}{\linewidth}{@{}Yl@{}}
    \toprule
    参数 & 参考值或设置\\
    \midrule
    感知节点数／接收中心数 & 9／1\\
    每任务候选上限 $K_{\max}$ & 4\\
    标称带宽 & 20 kHz\\
    数据／反馈发射功率 & 0.2 W／0.1 W\\
    参考信道增益 $\rho_0$ & $10^{-4}$\\
    噪声功率谱密度 $N_0$ & $4\times10^{-20}$ W/Hz\\
    每 token 字节数 $b$ & 4\\
    数据头部／有效载荷上限 & 64 B／1200 B\\
    反馈帧长度 & 24 B\\
    发送时长裕量 $\kappa$ & 1.6\\
    单次发送时长上限 $\tau_{\max}$ & 8 ms\\
    任务时限 $D_k$ & 25 ms\\
    反馈平均额外时延 $\bar d_{\mathrm f}$ & 6 ms\\
    反馈丢失率 $\epsilon_{\mathrm f}$ & 25\%\\
    速率波动强度 $\sigma_r$ & 0.55\\
    每节点初始发送能量 $E_i^{(0)}$ & 0.15 J\\
    数据尝试上限／基础等待间隔 & 9 次／2 ms\\
    C-BFC 能量／时间权重 & 0.24／0\\
    \bottomrule
  \end{tabularx}
\end{table}

\subsection{对比方法与评价指标}
\label{rq3:subsec:baselines}

经典对照采用固定冗余、停等 ARQ、选择重传、语义阈值 ARQ 和队列贪心五种方法。固定冗余按验证集选定的数量和副本数发送，停等 ARQ 等待确认或超时后继续，选择重传则允许窗口内未确认候选独立补发。语义阈值 ARQ 在已确认效用达到阈值时停止；队列贪心结合接收信念与能量价格，选择即时净收益最高的发送。

用于离散动作选择的稳定 Soft Actor-Critic 方法（Stable Discrete SAC，SD-SAC）通过熵惩罚和价值估计修正改善离散动作学习，本实验保留其双 Q 平均、Q 值裁剪和熵变化惩罚\citep{zhou2024sdsac}。两阶段深度 Q 网络控制器（Two-Phase Deep Q-Network，TP-DQN）参考两阶段码率控制，先预测语义效用，再以 Double DQN 训练动作价值\citep{wu2026twophase}。两者适配相同动作空间：四个发送、三个等待及停止；输入为候选、接收概率、效用和资源特征，网络为 $89\rightarrow128\rightarrow64\rightarrow8$。

两个控制器均在本环境重新训练，每个种子执行 80 个各含 600 任务的任务流，覆盖十二种物理与反馈设置。每四次交互更新，批量 128，学习率 $3\times10^{-4}$。SD-SAC 使用训练任务的正确完成奖励，TP-DQN 使用预测效用，二者均加入效用势函数差；测试参数由验证选择。两种对照均使用本章的候选语义表示和通信环境，分别保留其策略更新与两阶段控制机制。

五组条件分别为：速率波动 0.20、0.35、0.55、0.80、1.05；反馈额外时延 0、3、6、9、12 ms；任务时限 10、15、25、35、45 ms；每节点电量 0.06、0.10、0.15、0.24、0.36 J；组件实验反馈丢失率 0\%、15\%、25\%、45\%、60\%。每组仅改变对应参数。前四组去除共同参考点后为十七种链路与资源配置，含组件实验共二十一种配置，累计 1,089 次任务流回放。

主图展示固定冗余、语义阈值 ARQ、队列贪心、SD-SAC、TP-DQN 和 C-BFC，表中补充停等 ARQ 与选择重传。语义对照共享同种子的原始效用网络，C-BFC 另加一致性投影；阈值、重传规则及资源权重均经验证后固定。四联图的子图依次表示按时正确完成率、单位正确任务通信能耗、重复发送比例和平均通信完成时延；除各组所考察的参数外，其余条件保持参考设置。

实验从任务完成、能量支出、重复通信和时延四方面评价交付性能。

\paragraph{按时正确完成率}
截止前至少接收一条证据且分类正确计为完成：
\begin{equation}
  R_{\mathrm c}=\frac{1}{N_{\mathrm{task}}}\sum_{k=1}^{N_{\mathrm{task}}} u_k(S_k^{\mathrm{end}}).
  \label{rq3:eq:metric-correct}
\end{equation}
分母包含能量不足或直接停止的未服务任务。

\paragraph{单位正确任务通信能耗}
将整个任务流的数据与反馈发送能量除以正确完成的任务数：
\begin{equation}
  E_{\mathrm c}=\frac{\sum_{k=1}^{N_{\mathrm{task}}} E_k^{\mathrm{comm}}}
                         {\sum_{k=1}^{N_{\mathrm{task}}} u_k(S_k^{\mathrm{end}})}.
  \label{rq3:eq:metric-energy}
\end{equation}
单位为 mJ，先逐次回放计算，再跨重复实验及条件平均。

\paragraph{重复发送比例}
令 $d_{k,a}$ 表示发送开始前该证据是否已成功到达，则
\begin{equation}
  R_{\mathrm{dup}}=\frac{\sum_k\sum_{a\in\mathcal T_k}d_{k,a}}
                         {\sum_k|\mathcal T_k|}.
  \label{rq3:eq:metric-duplicate}
\end{equation}
该指标由实际日志计算；此前从未成功的补发不计为重复。

\paragraph{平均通信完成时延}
记 $\mathcal K_+=\{k:|\mathcal T_k|>0\}$ 为实际发起数据发送的任务集合，则
\begin{equation}
  T_{\mathrm{comm}}=\frac{1}{|\mathcal K_+|}\sum_{k\in\mathcal K_+}T_k.
  \label{rq3:eq:metric-latency}
\end{equation}
时延包含发送与等待，排除未发送任务，须结合正确率解释。

\section{实验结果与分析}

\subsection{总体性能比较}
\label{rq3:subsec:overall-results}

表~\ref{rq3:tab:main-results} 汇总十七种链路与资源配置。C-BFC 正确完成率为 49.84\%，单位正确任务能耗为 3.615 mJ，较固定冗余和语义阈值 ARQ 分别提高正确率 9.35 和 6.45 个百分点。前者不能按实际交付调整副本，后者在确认未返回时仍依赖超时。

\begin{table}[t]
  \centering\small
  \setlength{\tabcolsep}{3.5pt}
  \bicaption{十七个主比较条件下的平均性能}{Average performance across the seventeen main evaluation conditions}
  \label{rq3:tab:main-results}
  \begin{tabular*}{\linewidth}{@{\extracolsep{\fill}}lrrrr@{}}
    \toprule
    方法 & \shortstack{正确完成率\\(\%)} & \shortstack{单位正确任务\\能耗 (mJ)} & \shortstack{重复发送\\(\%)} & \shortstack{通信时延\\(ms)}\\
    \midrule
    固定冗余 & 40.49 & 4.941 & 29.35 & 11.32\\
    停等 ARQ & 43.65 & 3.678 & 7.49 & 16.69\\
    选择重传 & 43.57 & 3.665 & 7.16 & 16.77\\
    语义阈值 ARQ & 43.39 & 4.090 & 6.71 & 17.63\\
    队列贪心 & 47.26 & 4.564 & 25.76 & 10.96\\
    SD-SAC & 48.97 & 4.334 & 15.67 & 17.76\\
    TP-DQN & 49.18 & 4.219 & 19.11 & 17.13\\
    \midrule
    \textbf{C-BFC} & \textbf{49.84} & \textbf{3.615} & \textbf{17.14} & \textbf{16.47}\\
    \bottomrule
  \end{tabular*}
  \par\vspace{4pt}{\footnotesize 各条件等权汇总，共享的参考条件只计入一次。加粗行表示本章方法。}
\end{table}


队列贪心正确完成率为 47.26\%，但单位正确任务能耗达 4.564 mJ、重复比例为 25.76\%。缺少主动等待时，即时收益为正的连续发送可能消耗本可通过确认节省的能量。

相较 SD-SAC 和 TP-DQN，C-BFC 正确率提高 0.87、0.66 个百分点，单位正确任务能耗降低 16.59\%、14.32\%，时延降低 7.29\%、3.87\%；正确率分别在十二、十四个条件下更高，能耗在全部十七个条件下更低。其主要收益是在相近正确率下节能。

\begin{table}[t]
  \centering\small
  \setlength{\tabcolsep}{3.5pt}
  \bicaption{参考条件下的性能比较}{Performance comparison under the reference condition}
  \label{rq3:tab:reference-results}
  \begin{tabular*}{\linewidth}{@{\extracolsep{\fill}}lrrrr@{}}
    \toprule
    方法 & \shortstack{正确完成率\\(\%)} & \shortstack{单位正确任务\\能耗 (mJ)} & \shortstack{重复发送\\(\%)} & \shortstack{通信时延\\(ms)}\\
    \midrule
    固定冗余 & 41.68 & 5.023 & 29.86 & 11.79\\
    停等 ARQ & 45.66 & 3.627 & 7.72 & 17.46\\
    选择重传 & 45.50 & 3.615 & 7.40 & 17.46\\
    语义阈值 ARQ & 46.32 & 3.948 & 7.07 & 18.16\\
    队列贪心 & 47.66 & 4.706 & 27.20 & 11.36\\
    SD-SAC & 49.95 & 4.356 & 15.66 & 18.24\\
    TP-DQN & 49.86 & 4.329 & 21.49 & 17.74\\
    \midrule
    \textbf{C-BFC} & \textbf{51.08} & \textbf{3.639} & \textbf{17.40} & \textbf{17.56}\\
    \bottomrule
  \end{tabular*}
  \par\vspace{4pt}{\footnotesize 参考条件为 25 ms 任务时限、6 ms 反馈延迟、25\% 反馈丢失率及每节点 0.15 J 发送能量。}
\end{table}


参考条件见表~\ref{rq3:tab:reference-results}。C-BFC 正确率为 51.08\%，较学习对照高约 1.1 个百分点，单位正确任务能耗由约 4.3 mJ 降至 3.639 mJ。停等与选择重传虽重复较少，正确率仅约 45.5\%。等待确认能够避免部分重复发送，但也会占用补充有效证据的时间，因而不能仅凭重复比例判断交付效果。

\subsection{信道波动对交付效果的影响}
\label{rq3:subsec:fading-results}

图~\ref{rq3:fig:fading} 给出了不同速率波动强度下的结果。波动增大后，实际速率更容易偏离预测值，发送难以在预留时间内完成。C-BFC 正确率从 $\sigma_r=0.20$ 时的 57.26\% 降至 $\sigma_r=1.05$ 时的 41.32\%；在固定截止时间下，控制器用于补偿失败发送的机会也受到限制。

\begin{figure}[!htb]
  \centering
  \resultfigure{01_fading_compact.pdf}{不同速率波动强度下的性能比较}
  \bicaption{速率波动强度变化下的性能比较}{Performance under different rate fluctuation levels}
  \label{rq3:fig:fading}
\end{figure}

在 $\sigma_r=0.20$ 时，C-BFC、SD-SAC、TP-DQN 正确率均约 57\%，单位正确任务能耗分别为 2.346、3.464、2.831 mJ。重传需求较低时，接收信念和一致性约束有助于及早排除低收益发送。

在 $\sigma_r=1.05$ 时，三者能耗分别为 5.036、5.621、5.610 mJ，正确率差距缩小。与低波动条件相比，三种方法均需为一次正确完成付出更多通信能量，C-BFC 的节能幅度也随之减小。

\subsection{确认反馈延迟的影响}
\label{rq3:subsec:feedback-results}

图~\ref{rq3:fig:feedback} 比较不同反馈延迟下的交付性能。额外时延为零时，确认仍需串行传送；随着额外时延增加，控制器在较长时间内无法确定接收状态，确认返回后可用于调整发送的时间也相应缩短。

\begin{figure}[!htb]
  \centering
  \resultfigure{02_feedback_compact.pdf}{不同反馈延迟下的性能比较}
  \bicaption{反馈平均额外时延变化下的性能比较}{Performance under different additional feedback delays}
  \label{rq3:fig:feedback}
\end{figure}

额外时延为 0 ms 时，C-BFC、SD-SAC、TP-DQN 的正确率为 53.52\%、50.95\%、52.01\%，完成时延为 11.36、15.60、14.42 ms。及时确认使控制器更早停止重复发送或补充互补证据。

额外时延为 12 ms 时，三者正确率降至 48.92\%、48.10\%、48.52\%，单位正确任务能耗为 3.973、4.426、4.526 mJ。C-BFC 能耗较零延迟时的 3.447 mJ 上升，重复比例也有所增加。在给定任务窗口内，晚到确认留下的调整时间较少，主动等待所能节省的发送因而减少。

\subsection{任务截止时间的影响}
\label{rq3:subsec:deadline-results}

图~\ref{rq3:fig:deadline} 中，10 ms 时限往往只容纳一次主要发送，难以利用返回确认继续调整。C-BFC 与 SD-SAC 正确率分别为 44.48\% 和 44.64\%，差异很小。

\begin{figure}[!htb]
  \centering
  \resultfigure{03_deadline_compact.pdf}{不同任务截止时间下的性能比较}
  \bicaption{任务截止时间变化下的性能比较}{Performance under different task deadlines}
  \label{rq3:fig:deadline}
\end{figure}

时限为 45 ms 时，C-BFC 正确率达 54.74\%，较 SD-SAC 的 51.34\% 高 3.40 个百分点；单位正确任务能耗为 3.539 mJ，较其 4.466 mJ 低 20.76\%，两项指标也优于 TP-DQN。

该条件下 C-BFC 平均完成时延为 23.39 ms，低于 SD-SAC 的 28.20 ms 和 TP-DQN 的 26.03 ms。足够的任务窗口允许利用确认补发，并在后续收益不足时提前结束。

\subsection{节点发送能量的影响}
\label{rq3:subsec:energy-results}

图~\ref{rq3:fig:energy} 考察节点电量对连续任务的影响。每节点 0.06 J 时，C-BFC 正确率为 36.32\%，单位正确任务能耗为 2.427 mJ。由于任务之间不补充电量，当前任务多发送一条消息，就会减少后续任务可用的能量；低电量下这一约束更为明显。

\begin{figure}[!htb]
  \centering
  \resultfigure{04_energy_compact.pdf}{不同节点发送能量下的性能比较}
  \bicaption{每节点初始发送能量变化下的性能比较}{Performance under different initial node transmission energies}
  \label{rq3:fig:energy}
\end{figure}

电量增至 0.10 J 时，C-BFC 正确率为 46.03\%，高于 SD-SAC 的 44.66\%；在 0.15 J 时达 51.08\%。各条件采用同一组控制权重，发送价格则随剩余电量和虚拟队列变化，使策略能够根据实际支出调整后续发送。

每节点 0.36 J 时，C-BFC 正确率为 55.74\%，低于 SD-SAC 的 57.57\% 和 TP-DQN 的 56.76\%，但单位正确任务能耗为 4.608 mJ，低于二者的 6.164、6.033 mJ。高能量条件下仍保留节能偏好，因而放弃部分代价较高的交付机会。

\subsection{消融实验}
\label{rq3:subsec:ablation}

设置三个组件对照：无接收信念只使用最新位图，不传播未确认发送的成功概率，也不计算相应等待价值；无主动等待保留发送前瞻，但禁止主动等待；无一致性约束使用原始效用。各版本分别验证资源权重，再测试五档反馈丢失率，结果见图~\ref{rq3:fig:ablation} 与表~\ref{rq3:tab:ablation-results}。

\begin{figure}[!htb]
  \centering
  \resultfigure{05_ablation_compact.pdf}{不同反馈丢失率下的组件比较}
  \bicaption{不同反馈丢失率下的组件比较}{Component comparison under different feedback loss rates}
  \label{rq3:fig:ablation}
\end{figure}

\begin{table}[t]
  \centering\small
  \setlength{\tabcolsep}{3.5pt}
  \bicaption{五档反馈丢失率下的组件比较}{Component comparison at five feedback loss rates}
  \label{rq3:tab:ablation-results}
  \begin{tabular*}{\linewidth}{@{\extracolsep{\fill}}lrrrr@{}}
    \toprule
    方法 & \shortstack{正确完成率\\(\%)} & \shortstack{单位正确任务\\能耗 (mJ)} & \shortstack{重复发送\\(\%)} & \shortstack{通信时延\\(ms)}\\
    \midrule
    无接收信念 & 39.69 & 5.814 & 60.46 & 16.91\\
    无主动等待 & 48.89 & 4.516 & 32.59 & 10.96\\
    无一致性约束 & 51.28 & 4.064 & 16.96 & 18.63\\
    \midrule
    \textbf{C-BFC} & \textbf{50.97} & \textbf{3.674} & \textbf{18.29} & \textbf{17.62}\\
    \bottomrule
  \end{tabular*}
  \par\vspace{4pt}{\footnotesize 各版本分别完成验证参数选择；表中数值为五个反馈丢失条件的等权平均。}
\end{table}


去除接收信念使平均正确率由 50.97\% 降至 39.69\%，单位正确任务能耗由 3.674 mJ 增至 5.814 mJ，重复比例达 60.46\%。仅依赖已确认位图会在反馈缺失时低估已有证据，导致重复分配资源。

去除等待后正确率为 48.89\%，单位正确任务能耗为 4.516 mJ，重复比例为 32.59\%。完整方法以完成时延从 10.96 ms 增至 17.62 ms 为代价，换取更准确的状态信息和更少的连续发送。

无一致性约束的正确率为 51.28\%，略高于完整方法的 50.97\%，能耗则为 4.064 mJ。投影将其降至 3.674 mJ，降幅约 9.6\%，伴随 0.31 个百分点的正确率变化；作用主要体现为抑制虚假效用增量、减少通信支出。


\section{本章小结}
\label{rq3:subsec:discussion}

本章提出 C-BFC，根据发送记录与累计确认估计接收状态，并比较发送、等待和停止对任务完成的作用。同一任务中，支持相同判断的消息组合采用一致的收益估计，以减少不改变判断结果的重复发送；累计能耗用于调节后续任务的支出。十七种链路与资源配置下，平均正确完成率为 49.84\%，单位正确任务能耗为 3.615 mJ，较 SD-SAC 和 TP-DQN 分别降低 16.59\% 和 14.32\%。方法在相近交付水平下减少通信支出，收益仍受时限与能量条件制约。
